% Related work section for OrbitGraph paper.

\section{Related Work}
\label{sec:related}

We group prior work by routing strategy, SDN management in satellite networks,
evaluation platforms, and broader LEO networking literature, and contrast each
line with OrbitGraph's graph-centric, proactive SDN forwarding model
(Section~\ref{sec:orbitgraph}).

\subsection{Centralized and snapshot routing}

Chang et al.~\cite{seong:topological-design:1995} model a LEO constellation as
a finite-state automaton whose states correspond to recurring visibility
configurations.
For each state they jointly optimize inter-satellite link assignment and
multi-commodity routing offline via simulated annealing and linear programming,
then preload tables on board.
This snapshot strategy avoids transient instability during topology transitions
but relies on centralized offline computation, assumes moderate constellation
size (12 satellites in the reported experiments), and does not address online
failure recovery.
OrbitGraph shares the \emph{snapshot} view---each emulation second yields a graph
$G_t$---but computes shortest paths online and pushes kernel routes through an
SDN controller rather than storing per-state table sequences on satellites.

Handley~\cite{handley:ground-relays:2019} studies wide-area routing in
mega-constellations with and without inter-satellite links, emphasizing
ground-relay paths when ISLs are absent and hierarchical, tile-based updates
when centralized Dijkstra over thousands of relays would be too slow.
The work is centralized and delay-oriented like OrbitGraph, but targets
constellation \emph{design} and ground-side relay selection rather than
in-container forwarding during live GSL handover; it does not install routes on
emulated satellite nodes or compare against a distributed IGP baseline.

\subsection{Distributed LEO routing protocols}

Stock et al.~\cite{gregory:satellite-routing:2022} propose DisCoRoute, a
distributed on-demand algorithm for Walker-Delta mega-constellations.
They derive a closed-form minimum hop count between orbital identifiers and
select the next hop in constant time without global Dijkstra, reporting large
speedups over shortest-path search at scales up to 32\,000 satellites.
DisCoRoute optimizes hop count, assumes a fixed +Grid ISL pattern, and does
not model link failures or data-plane convergence delay.
OrbitGraph instead centralizes shortest-path computation on delay-weighted
graphs and measures end-to-end outage during handover; the two approaches are
complementary (distributed hop logic vs.\ timed SDN install).

Pan et al.~\cite{pam:opspf:2019} present OPSPF, which keeps OSPF-style hello
and link-state update machinery but uses orbit prediction so that regular
constellation motion generates no protocol traffic; only irregular events
(failures) trigger updates.
Evaluated on 48 satellites, OPSPF reduces convergence time relative to vanilla
OSPF.
OrbitGraph removes distributed flooding entirely for the SDN mode: the
controller recomputes $G_t$ from ephemeris and programs routes directly.
OPSPF remains a distributed-resilience design; OrbitGraph targets zero handover
outage via proactive install ordering rather than faster LSU propagation.

Our OSPF baseline (BIRD on StarryNet~\cite{lai2023starrynet}) represents the
reactive link-state approach analyzed in early protocol
studies~\cite{sidhu:ospf:1993} and widely used in academic LEO
evaluations~\cite{manzanareslopez2025comprehensive}.
It reconverges only after full topology mutation, which Section~\ref{sec:analysis}
shows can black-hole traffic for multiple seconds at GSL handover---the behavior
OrbitGraph is designed to eliminate.

Terrestrial ad hoc routing (AODV~\cite{perkins:aodv:1999},
GPSR~\cite{karp:gpsr:2000}) targets unpredictable mobility without orbital
models; ns-3-leo reports that such MANET protocols perform poorly in highly
dynamic satellite scenarios~\cite{ns3leo2022}.
OrbitGraph exploits the opposite assumption: topology is predictable from
SGP4-based delay matrices, so a centralized graph controller is feasible.

\subsection{SDN and graph-based network management}

The SDN paradigm separates control and data planes and enables programmatic
forwarding through interfaces such as OpenFlow~\cite{10.1145/1355734.1355746}.
Pantuza et al.~\cite{7014202} show that representing an SDN domain as a graph
gives controllers a unified structure for topology-aware management and
shortest-path algorithms; OrbitGraph applies that graph-centric control model
to time-varying LEO snapshots rather than terrestrial OpenFlow switches.

Satellite-specific SDN research has focused on \emph{controller placement},
not data-plane route computation.
Papa et al.~\cite{papa:sdn-cpp:2018} formulate dynamic controller placement on
an Iridium-class 66-satellite constellation as an integer linear program that
minimizes average flow setup time under time-varying traffic.
Guo et al.~\cite{guo:spda:2022} propose static placement with dynamic
assignment (SPDA): controllers remain on fixed satellites while switch-to-
controller bindings adapt to topology snapshots, reducing migration cost on
Iridium and Celestri models.
Both lines optimize \emph{where} control logic runs and \emph{which} nodes attach
to which controller; they treat forwarding as a shortest-path black box.
OrbitGraph addresses the complementary problem---\emph{what} routes to install
and \emph{when} relative to link lifetime---and evaluates data-plane availability
during handover, not switch-controller latency alone.

\subsection{Simulation and emulation platforms}

Several tools support LEO routing research at different fidelity levels.
The ns-3-leo module~\cite{ns3leo2022} extends discrete-event simulation with
orbital mobility, ISLs, and MANET routing protocols, highlighting limitations
of terrestrial protocols in space.
Kempton and Riedl~\cite{kempton:network-simulator:2021} provide a dedicated
LEO network simulator with configurable +Grid and sparse ISL topologies and
3D visualization; the underlying SILLEO-SCNS tool~\cite{kempton2020silleo,
online:silleo:2025} focuses on topology generation rather than packet-level
stacks.
Dewangan and Sahu~\cite{dewangan:simple-omnet:2017} demonstrate simplified
satellite scenarios in OMNeT++~5 for rapid prototyping.
These platforms either lack real protocol stacks (SILLEO-SCNS) or use
distributed routing without centralized proactive install (ns-3-leo with AODV).

Packet-level ns-3 frameworks are surveyed by
Manzanares-Lopez et al.~\cite{manzanareslopez2025comprehensive}.
Hypatia~\cite{kassing2020exploring} precomputes orbital state and runs ns-3
with static Floyd--Warshall routing; StarPerf~\cite{lai2020starperf} profiles
area-to-area performance statistically; LEOPath~\cite{leopath2024} targets
routing-algorithm comparison without end-to-end kernel forwarding.
StarryNet~\cite{lai2023starrynet} emulates real TCP/IP and OSPF in containers
and is our evaluation harness; OrbitGraph extends it with a centralized
controller, incremental kernel route install, and proactive handover
(Section~\ref{sec:simulation}).

\subsection{Broader LEO networking literature}

Growing interest in large LEO and swarm systems motivates global, low-latency
connectivity~\cite{smallsats2020,asri2024swarm,farrag:swarm-survey:2021}.
Swarm intelligence networking frameworks~\cite{chen2022swarm} and formation
control studies~\cite{fangchen2022pdg} address coordination at the mission
layer rather than IP handover semantics; they do not replace routing protocols
but illustrate the scale and dynamics OrbitGraph targets at the network layer.

Zhu et al.~\cite{zhu:ml-routing-survey:2025} survey machine-learning routing
for LEO networks, covering model design, training, and deployment challenges.
That literature pursues learned policies; OrbitGraph uses explicit graph
snapshots and Dijkstra shortest paths with deterministic SDN install ordering,
which yields reproducible handover behavior suitable for controlled comparison
with OSPF.

Our earlier experimentation environment for distributed routing analysis in LEO
swarms~\cite{online:silleo-routing:2025} compared conventional protocols without
a centralized proactive strategy.
The present work advances that line by implementing OrbitGraph---graph-centric
SDN with incremental, make-before-break updates---and quantifying data-plane
outage against the same class of emulated constellations.

\subsection{Positioning of OrbitGraph}

Table~\ref{tab:related-positioning} summarizes how OrbitGraph relates to the
closest prior lines.
Prior snapshot and centralized schemes optimize offline tables or relay
placement; distributed protocols optimize hop count or OSPF convergence;
satellite SDN papers optimize controller placement.
OrbitGraph combines online graph snapshots, centralized delay-weighted
shortest paths, incremental dataplane deltas, and proactive GSL handover in a
single emulated system with measured end-to-end outage---a combination not
present in the works above.

\begin{table}[!t]
  \centering
  \caption{OrbitGraph vs.\ representative related work (selected dimensions).}
  \label{tab:related-positioning}
  \renewcommand{\arraystretch}{1.1}
  \footnotesize
  \begin{tabular}{@{}lcccc@{}}
    \hline
    Work & Central & Real FIB & Proactive HO & Outage eval. \\
    \hline
    Chang et al.~\cite{seong:topological-design:1995} & Yes & No  & offline & No \\
    DisCoRoute~\cite{gregory:satellite-routing:2022} & No  & No  & No      & No \\
    OPSPF~\cite{pam:opspf:2019} & No  & Yes & partial & partial \\
    Papa / Guo SDN~\cite{papa:sdn-cpp:2018,guo:spda:2022} & Yes & No  & No      & No \\
    StarryNet OSPF~\cite{lai2023starrynet} & No  & Yes & No      & Yes \\
    \textbf{OrbitGraph (this work)} & Yes & Yes & Yes     & Yes \\
    \hline
  \end{tabular}
  \\[0.4em]
  \footnotesize
  HO: handover. ``Real FIB'': kernel or protocol forwarding on emulated nodes.
\end{table}
